% Written By Enoch Yu
% Downloaded from archive.enochyu.com

\documentclass{article}

\usepackage{amsmath}
\usepackage{amssymb}
\usepackage{amsthm}
\usepackage{mathtools}

\usepackage[margin=1in, headsep=15pt]{geometry}
\usepackage{lastpage}
\usepackage{fancyhdr}
\fancyhead[R]{Enoch Yu}
\fancyfoot[C]{Page \thepage \hspace{1pt} of \pageref{LastPage}}
\pagestyle{fancy}
\usepackage{parskip}

\usepackage{enumitem}
\usepackage[dvipsnames]{xcolor}
\usepackage[normalem]{ulem}
\usepackage{url}
\usepackage{hyperref}
\hypersetup{
    colorlinks=true,
    linkcolor=black,
    filecolor=magenta,      
    urlcolor=cyan,
    pdftitle={7.8.2025 AMC 10 Free Class Preview Homework Solution},
    pdfpagemode=FullScreen,
}
\usepackage{tabularx}
\usepackage{arydshln}
\usepackage{float}
\usepackage{pdfpages}

\usepackage{tikz}
\usetikzlibrary{calc, angles, shapes.geometric}
\usepackage{tkz-euclide}
\usepackage{pgfplots}
\pgfplotsset{compat=1.9}
\usepackage[font=small,labelfont=bf]{caption}

\newtheorem{theorem}{Theorem}[section]
\newtheorem{lemma}[theorem]{Lemma}
\newtheorem*{lemma*}{Lemma}
\newtheorem{sublemma}{Lemma}[section]
\newtheorem{proposition}{Proposition}
\newtheorem{corollary}{Corollary}[theorem]
\newtheorem{example}{Example}[section]
\newtheorem*{example*}{Example}
\newenvironment{solution}{\begin{trivlist}\item[]{\bf Solution}}{\qed \end{trivlist}}

\title{7.8.2025 AMC 10 Free Class Preview Homework Solution}
\author{Enoch Yu}
\date{July 2025}

\begin{document}

\maketitle

\subsection*{Problem}
Find the number of real solutions of the equation
$x|x|- 3|x| + 2 = 0$.
\begin{solution}
First and foremost, the absolute values could be removed.
\[
\begin{cases}
    x \ge 0 \implies x^2 - 3x + 2 = 0 \\
    x < 0 \implies -x^2 + 3x + 2 = 0
\end{cases}
\]
For the first case, the solutions to the equation are
$1$ and $2$. The roots are valid since both $1$ and $2$
are greater than zero. The solutions to the second equation
are $\frac{3 \pm \sqrt{17}}{2}$. The only valid solution
from the second equation is $\frac{3 - \sqrt{17}}{2}$
because the root must be negative. Therefore, the number
of real solutions to the equation is $\boxed{3}$.
\end{solution}

\subsection*{Problem}
Find the number of pairs of non-negative integer solutions
to the equation $|a - b| + ab = 1$.
\begin{solution}
The absolute value could be removed.
\[
\begin{cases}
    a \ge b \implies ab + a - b = 1 \\
    a < b \implies ab - a + b = 1
\end{cases}
\]
The first equation could be factored using Simon's Favorite
Factoring Trick. From $(a - 1)(b + 1) = 0$, the possible
ordered pairs of $(a, b)$ are $(1, 0)$ and $(1, 1)$.

The second equation will lead to $(a + 1)(b - 1) = 0$.
The possible ordered pair is $(0, 1)$. Therefore, there
exists $\boxed{3}$ ordered pairs of non-negative integer
solutions to the equation $|a - b| + ab = 1$.
\end{solution}

\subsection*{Problem}
What is the product of all the roots of the equation
$\sqrt{5|x| + 8} = \sqrt{x^2 - 16}$?

\begin{solution}
Because square root functions are one to one correspondence,
both sides could be squared.
\[
5|x| + 8 = x^2 - 16
\]
Therefore, the absolute value could be removed.
\[
\begin{cases}
    x \ge 0 \implies x^2 - 5x - 24 = 0 \\
    x < 0 \implies x^2 + 5x - 24 = 0
\end{cases}
\]
The first equation gives $x = 8$. The second equation
provides $x = -8$. Moreover, notice that if $x = \pm 8$,
the left hand side and the right hand side will lead to a
valid real equalities. Therefore, the product of all the
roots of the equation is $\boxed{-64}$.
\end{solution}

\subsection*{Problem}
What is the product of the real roots of the equation
$x^2 + 18x + 30 = 2\sqrt{x^2 + 18x + 45}$?

\begin{solution}
Let $A = x^2 + 18x + 30$. In other words,
$A = 2\sqrt{A + 15}$. Therefore, $A^2 - 4A - 60 = 0$.
However, $A$ can only be $10$ since $A$ must be a
positive number. $x^2 + 18x + 10 = 0$ will lead to real
roots with the product of $\boxed{-10}$.
\end{solution}

\subsection*{Problem}
Given that $x^2 + y^2 = 10$, $\sqrt[4]{xy} + \sqrt{xy}
+ 27 = 29$, $x > 0$ and $y > 0$. What is $x+y$?
\begin{solution}
Let $a = \sqrt[4]{xy}$. Therefore, $a^2 + a - 2 = 0$.
In other words, $\sqrt[4]{xy} = -2, 1$. Because $x$ and
$y$ are real numbers, $xy$ must be $1$. Thus
$(x + y)^2 = 10 + 2$. In other words, $x + y = \boxed{2\sqrt{3}}$.
\end{solution}
\end{document}

